\section{Fourier structure of the tree code}\label{sec:tree-fourier}

The cut depends on the generated word, so every transformation that
preserves that word also constrains the cut's row Fourier coefficients.
At each gate these constraints give a common Fourier index shared by
all child pairs. Negating the gate table distinguishes whether this
index is an evaluation label. We then show that, at a random slice,
little smoothed Fourier mass is lost by keeping the entire common index
in just one child of each pair. Finally we express these retained
indices in row coordinates that remain free after the other rows are
fixed. This is the local form used in \Cref{sec:extraction}.
Throughout the section, $f$ is a function of the code in the sense of
\eqref{eq:function-of-code}.

\subsection{Row Fourier support}

For each node $v$, define its \emph{generator} $P(v)$, a set of leaves,
recursively: $P(v)=\{v\}$ at a leaf, and
$P(v)=P(v_{1,0})\cup P(v_{1,1})$ at an internal node.
Thus a generator always follows both children of the first pair;
unlike the slice, it involves no random choices.
Adding the same $h\in H$ to all rows in $P(v)$ multiplies the
subtree output at $a$ by $(-1)^{h(a)}$.  This follows by induction
from Equation~\eqref{eq:gate-rule}: multiplying both children of
one pair by the same sign leaves all gate inputs $z_g$ unchanged
and multiplies the prefactor by that sign.

For a row Fourier index $C=(q_i)_{i=1}^k\in(H^*)^k$, put
\[
 \chi_C(R)=\prod_i\chi_{q_i}(R_i),\qquad
 |C|=|\{i:q_i\ne0\}|,\qquad
 q_v=\sum_{i\in P(v)}q_i.
\]
Write $q_{v,gb}$ for the index of child $(g,b)$ of $v$.
Thus, by definition, $q_v=q_{v,10}+q_{v,11}$.
At fixed tables, the row expansion and its smoothing are
\begin{equation}
 f(R,F)=\sum_C\widehat f_C(F)\chi_C(R),\qquad
 \widehat{\bar f}_C(F)=r^{|C|}\widehat f_C(F).
 \label{eq:row-smoothing}
\end{equation}

\begin{lemma}[Row symmetries]
\label{lem:row-symmetry}
Let $f$ be a function of the code.  Every nonzero row Fourier
coefficient of $f$ satisfies, at each internal node $v$,
\begin{equation}
 q_{v,g0}+q_{v,g1}=q_v\qquad(g\in[s]).
 \label{eq:row-relations}
\end{equation}
For a table collection $F$, let $F^{v,-}$ negate $F_v$ alone,
and let $F^{v,z}$ replace $F_v(u)$ by $F_v(u+z)$, leaving all
other tables unchanged.  Then
\begin{align}
 \widehat f_C(F^{v,-})
   &=(-1)^{q_v(\one)}\widehat f_C(F),
   \label{eq:table-flip-action}\\
 \widehat f_C(F^{v,z})
   &=(-1)^{\sum_g z_g q_{v,g1}(\one)}\widehat f_C(F).
   \label{eq:table-translation-action}
\end{align}
All these statements remain true after row smoothing.
They are also preserved by row Fourier projections and by
projection to any range of total Fourier degrees in the entries
of a specified collection of whole gate tables.
\end{lemma}

\begin{proof}
For each $g$, shifting the generators of both children in pair
$g$ by $h$ has exactly the same effect on the subtree at $v$ as
shifting $P(v)$ by $h$.  Consequently the two row translations
have the same effect on the whole word.  Taking row Fourier
coefficients yields
\[
 \chi_{q_{v,g0}+q_{v,g1}}(h)\widehat f_C(F)
   =\chi_{q_v}(h)\widehat f_C(F)\qquad(h\in H).
\]
Distinct characters are distinct functions, proving
Equation~\eqref{eq:row-relations}.

Negating $F_v$ has the same effect on the word as adding
$\one$ to all rows in $P(v)$.  This proves
Equation~\eqref{eq:table-flip-action}.  To translate the input of
$F_v$ by $z$, add $\one$ to the generator of its second child in
each pair with $z_g=1$.  This changes precisely those gate input
bits and leaves its prefactor unchanged.  Fourier transformation
gives Equation~\eqref{eq:table-translation-action}.

Row smoothing and row Fourier projections act diagonally on
$C$, so preserve all three identities.  On the cube of
table-entry signs, a domain translation permutes coordinates,
and negating a gate table multiplies each monomial by its
parity sign.  Both operations preserve total degree in any
specified collection of whole gate tables and therefore
commute with the stated degree projections.
\end{proof}

\subsection{An unsplit projection on a random slice}

Equation~\eqref{eq:row-relations} assigns one common sum $q_v$ to every
child pair. Equation~\eqref{eq:table-flip-action} then has a useful
interpretation: coefficients odd under negating $F_v$ have
$q_v(\one)=1$, so their common index is a label by
\eqref{eq:labels-dual}. The individual child indices, however, can
still be arbitrary functionals whose sum is $q_v$.

We simplify this last freedom at one gate at a time. Keeping $q_v$ in
one child of each pair leaves only the choice of that child, rather
than two arbitrary functionals with prescribed sum.
Call $C$ \emph{unsplit at $v$} if, for every $g$, the ordered pair
$(q_{v,g0},q_{v,g1})$ is either $(q_v,0)$ or $(0,q_v)$.
Let $S_v$ be the orthogonal row Fourier projection onto such
indices.  When $q_v=0$, this requires every child index to be zero.
The next bound averages the cost over the slice depth. It does not
require an index to be unsplit at every gate of the tree.

\begin{lemma}[Average cost of the unsplit projection]
\label{lem:unsplit}
Let $f$ be a function of the code with $|f|\le1$, let
$r=1-\mu/d$ with $0<\mu<1/2$, and let $\bar f=T_rf$.
Choose the random slice independently of $f$.  Then
\begin{equation}
 \E_{J_0,\mathcal D}\sum_{v\in\mathcal D}
       \norm{(I-S_v)\bar f}_{2,R,F}^{2}
       \le \frac{s}{2e m\mu}.
 \label{eq:unsplit-bound}
\end{equation}
The same bound holds after averaging over additional background
data determining $f$, provided the slice is chosen independently
of those data.
\end{lemma}

\begin{proof}
Fix an index $C$ satisfying Equation~\eqref{eq:row-relations},
and let $N_j$ be the number of nodes at depth $j$ with nonzero
index $q_v$.  If $q_v\ne0$, each child pair contains at least
one nonzero index.  A failure of the unsplit condition at such a
node supplies at least one additional nonzero child.  If $q_v=0$,
a failure supplies at least two nonzero children, whereas its
minimum contribution is zero.  Therefore the number $B_j$ of
failures at depth $j$ satisfies
\[
 B_j\le N_{j+1}-sN_j.
\]
Using the slice inclusion probability and telescoping gives
\begin{align*}
 \E_{J_0,\mathcal D}\sum_{v\in\mathcal D}
          \one_{\{C\text{ is split at }v\}}
 &\le \frac1m\sum_{j=0}^{m-1}
          \frac{N_{j+1}-sN_j}{s^j}\\
 &=\frac{s}{m}\left(\frac{N_m}{s^m}-N_0\right)
 \le\frac{s|C|}{md}.
\end{align*}
Here $N_m=|C|$.  By Equation~\eqref{eq:row-smoothing} and
Parseval, the left side of Equation~\eqref{eq:unsplit-bound}
is at most
\[
 \frac{s}{m}\sum_C
   \left(r^{2|C|}\frac{|C|}{d}\right)
           \E_F|\widehat f_C(F)|^2.
\]
For $x\ge0$, $xe^{-2\mu x}\le(2e\mu)^{-1}$, and
$\log r\le-\mu/d$.  Hence every parenthesized factor is at
most $(2e\mu)^{-1}$, while
$\sum_C\E_F|\widehat f_C(F)|^2=\norm f_2^2\le1$.
This proves the bound.  The proof is pointwise in the background
data, so its average satisfies the same inequality.
\end{proof}

\subsection{Free shifts and the remaining row coordinates}

The extraction argument must use the common index after fixing most
of the rows. We therefore choose coordinates in which translating a
child generator changes one free row and leaves all other coordinates
fixed. The unsplit projection will then act separately at each fixed
context, and its retained characters will have an explicit form.

Fix a slice $\mathcal D$ and a node $v\in\mathcal D$.
The $2s$ sets $P(v_{g,b})$ are nonempty and disjoint.
Choose one representative leaf $i_{g,b}$ from each of them,
using a rule depending only on the tree, and set
$L_{g,b}=R_{i_{g,b}}$.  For every other leaf $i$ in that
generator, replace $R_i$ by $R_i+L_{g,b}$.  Keep all rows outside
these generators as they are.  Denote the resulting $k-2s$
remaining coordinates by $\Gamma\in H^{k-2s}$.

\begin{lemma}[Free-shift coordinates]
\label{lem:shift-coordinates}
The change of coordinates
\[
 H^k\longrightarrow H^{2s}\times H^{k-2s},\qquad
 R\longmapsto((L_{g,b}),\Gamma)
\]
is a linear bijection.  In particular, under uniform rows the
free shifts are independent uniforms conditional on $\Gamma$.
The Fourier index of $L_{g,b}$ corresponding to a full row
index $C$ is $Q_{g,b}=q_{v,g b}$.

Consequently, at every fixed $\Gamma$, the operator $S_v$ acts
as the Fourier projection in the $2s$ shifts onto precisely the
unsplit shift indices.  For a function of the code or its row
smoothing, the shift coefficients vanish unless the sums
$Q_{g,0}+Q_{g,1}$ are the same for every $g$.  The table actions
in Equations~\eqref{eq:table-flip-action} and
\eqref{eq:table-translation-action} hold for each shift
coefficient at this fixed context, with
$q_v=Q_{1,0}+Q_{1,1}$ and $q_{v,g1}=Q_{g,1}$.
\end{lemma}

\begin{proof}
The inverse map sets the representative row equal to $L_{g,b}$
and recovers every other row in its generator by adding
$L_{g,b}$ to its recorded difference.  This proves bijectivity.
In a row character, the coefficient of $L_{g,b}$ is exactly
$\sum_{i\in P(v_{g,b})}q_i=q_{v,g b}$.
Thus the condition defining $S_v$ depends only on the shift
index $Q$, not on any Fourier index of $\Gamma$.  Grouping the
full Fourier expansion by $Q$ proves its asserted action on
each fixed fiber.  The same grouping proves the support and
table-action identities: all full row coefficients with a
given $Q$ satisfy the same identities and carry the same signs.
\end{proof}

Thus, for a nonzero common index $q$, the retained shift indices are
specified by the subset $I\subseteq[s]$ of pairs in which $q$ occurs
in the second child. Their characters are
\[
 \prod_{g\notin I}\chi_q(L_{g0})
 \prod_{g\in I}\chi_q(L_{g1}).
\]
If $q=0$, there is only the all-zero unsplit index, with character $1$.
Negating $F_v$ multiplies the corresponding coefficient by
$(-1)^{q(\one)}$. On the part odd in $F_v$, $q$ is therefore a label
and hence nonzero, so the description by $q$ and $I$ is unique.
For such a label, translating the argument of $F_v$ by $z\in\F^s$
multiplies its coefficient by $(-1)^{\sum_{g\in I}z_g}$.
Thus the same subset $I$ records both the row character and the table
translation character. This is the local Fourier description used
in \Cref{sec:extraction}, valid at each fixed context.
